\chapter{Motivation}
\label{chap:motivation}

\section{The Need for Contactless Physiological and Occupancy Sensing}
Human occupancy detection and vital sign tracking are foundational pillars for intelligent environments. In elderly care facilities, independent living apartments, and hospital wards, monitoring the presence and physiological state of individuals is critical for preventing unattended falls, detecting respiratory distress, and ensuring rapid emergency response. According to global healthcare statistics, respiratory irregularities are among the earliest indicators of physiological deterioration, cardiopulmonary disease, and sleep apnea. 

However, conventional clinical sensing systems rely heavily on contact-based apparatus, such as pulse oximeters, electrocardiogram (ECG) leads, and thoracic chest belts. In everyday living environments, these devices face severe patient compliance barriers. Individuals frequently forget to wear devices, find adhesive electrodes uncomfortable, or experience skin irritation. Similarly, smartwatches and fitness bands require daily battery recharging and conscious adherence. A truly pervasive sensing system must operate passively in the background without imposing any behavioral burden on the user.

\section{Privacy Preservation vs. Visual Surveillance}
Optical video cameras have traditionally dominated smart surveillance and activity monitoring. Despite advances in computer vision and deep learning, visual surveillance encounters vehement societal resistance when deployed in private spaces such as bedrooms, bathrooms, and assisted living quarters. The inherent risk of video feed interception, unauthorized data harvesting, and cloud exposure creates substantial privacy liabilities. Furthermore, optical cameras are vulnerable to line-of-sight obstructions, poor ambient illumination, and optical obscuration caused by smoke or steam.

Radio frequency (RF) sensing fundamentally resolves this privacy paradox. Wireless electromagnetic waves do not capture identifiable facial features, skin tones, or personal images. Instead, they interact with physical mass and dielectric boundaries, capturing only anonymous spatial kinematics and thoracic oscillations. By utilizing RF Channel State Information, occupancy and health status can be continuously assessed with complete visual anonymity.

\section{Economic Viability: Democratizing RF Sensing with \$4 Microcontrollers}
For over a decade, wireless physical-layer sensing remained locked inside academic research laboratories due to extreme hardware barriers. Seminal CSI frameworks such as the Linux 802.11n CSI Tool relied on obsolete Intel 5300 or Atheros AR9580 mini-PCIe wireless cards. These cards necessitated bulky desktop PCIe backplanes, patched Linux kernels, high electrical power consumption, and equipment costs exceeding \$200 to \$1500 per sensing station.

Recent advances in Internet of Things (IoT) semiconductors, exemplified by the Espressif ESP32 System-on-Chip (SoC), present an unprecedented opportunity to democratize RF sensing. The ESP32 integrates a dual-core 240~MHz 32-bit processor, Wi-Fi 802.11b/g/n transceiver, and hardware baseband accelerators on a single silicon die retailing for approximately \$4. By developing lightweight digital signal processing algorithms capable of executing on edge microcontrollers, sophisticated RF human sensing can be deployed ubiquitously in homes, hospitals, and educational institutions at a fraction of conventional costs.

\section{Bridging Theory to Practical Edge Deployment}
While recent academic literature has proposed complex deep learning architectures (e.g., CNN-LSTMs, Graph Neural Networks) for CSI classification, these models require powerful GPUs, high inference energy, and massive training datasets that overfit to specific room layouts. When transferred to a new room with different furniture arrangements, such models experience catastrophic performance drops. There exists a critical necessity for mathematically grounded, lightweight DSP pipelines that operate in real time without offline training, providing robust plug-and-play occupancy and vital sign intelligence at the network edge.
